
% ============================================================================
%  فهرست گزینه‌ها  -- همه‌ی کلیدهای زیر اختیاری‌اند.
%  هر خط را در بلوک \atlset{...} بالا کپی کنید (کلیدها با ویرگول انگلیسی جدا می‌شوند).
%  قالب نوشتن:   key = value,          واژه یا عدد ساده به آکولاد نیاز ندارد
%                key = {دو واژه، یا ویرگول},     متن دارای فاصله یا ویرگول در آکولاد بیاید
%  اولویت: پیش‌فرض‌های قالب < همین پرونده < preamble/local.tex  (مورد بعدی غالب است).
%  مسیرها نسبت به همین پوشه (latex/) نوشته می‌شوند، مثلاً  logo = fig/logo.png
%  نام کلیدها و مقدارهای انتخابی را لاتین و بدون آکولاد بنویسید.
% ----------------------------------------------------------------------------
%  ظاهر
%    palette    = blue | red | gray | cyan | magenta | yellow | green | sunset | ocean | berry | neon
%    listing    = darcula | dark | bright | github           (پوسته‌ی رنگی بلوک کد)
%    fafont     = vazirmatn | felmi                          (فونت متن فارسی؛ پیش‌فرض vazirmatn)
%    fascale    = 0.95                                       (مقیاس فونت فارسی، عدد اعشاری با نقطه؛ vazirmatn: 0.95، felmi: 1.1)
%    headfont   = sans | serif                               (سرتیترها و برچسب‌ها)
%    headsize   = normal | compact                           (اندازه‌ی سرتیتر؛ compact برای مقاله)
%  صفحه‌آرایی
%    margins    = normal | narrow | formal | wide            (formal: ۳ سانتی‌متر سمت صحافی، یعنی راست)
%    linespread = 1.15                                       (عدد اعشاری بدون یکا؛ ۱٫۵ برابر: 1.5)
%    parstyle   = block | indent                             (بند با فاصله / بند با تورفتگی)
%    numbering  = flat | section                             (section: شماره‌ی معادله، شکل و جدول به‌صورت ۲-۱، ۲-۲ ...)
%    sectionbreak = off | on                                 (on: هر \section از صفحه‌ی تازه شروع شود)
%  صفحه‌ی عنوان
%    titlepage  = block | cover | classic | none             (بنر / جلد کامل / سبک مقاله / هیچ)
%    infofields = {author,studentid,course,instructor,duedate}
%                 فهرست کلیدهایی که در جدول مشخصات می‌آیند؛ جداشده با ویرگول و بدون فاصله. کلیدهای مجاز:
%                 author studentid course coursecode instructor supervisor advisor team partners
%                 experiment labsection term date duedate submissiondate performed university department
%    doctype    = {گزارش نهایی}                              (جایگزین برچسب خودکار «تمرین»، «گزارش» و ...)
%    logo       = fig/logo.png                               (فقط جلد؛ png، jpg یا pdf)
%  مشخصات سند  (متن؛ مقدار تهی نادیده گرفته می‌شود)
%    title subtitle author studentid course coursecode instructor supervisor advisor team partners
%    experiment labsection term date duedate submissiondate performed university department
%    affiliation keywords
%    date = {\today}   (تاریخ شمسی)     یا    date = {۲۱ اسفند ۱۴۰۴}
%    keywords = {واژه‌ی نخست، واژه‌ی دوم}                      (مقاله و جستار: زیر چکیده)
%    affiliation = {گروه \\ دانشگاه}                          (صفحه‌ی عنوان classic؛ \\ خط را می‌شکند)
%  سرصفحه و مراجع
%    header      = on | off                                  (سرصفحه‌ی جاری)
%    headerstart = {عنوان کوتاه}                             (سمت آغاز، یعنی راست؛ پیش‌فرض: درس، وگرنه عنوان)
%    headerend   = {نام فصل}                                 (سمت پایان؛ پیش‌فرض: بخش جاری)
%    bibstyle    = IEEEtran                                  (هر سبک BibTeX نصب‌شده، مثل plain یا ieeetr)
%  تغییر در میانه‌ی نوشتن:   \atlsetpalette{green}   \atlsetlisting{bright}   (docs/USER-MANUAL.md)
% ============================================================================
